\section{Introduction}
\label{sec:introduction}

Late-interaction retrieval models like ColBERT \citep{khattab2020colbert} achieve 
strong ranking quality by representing queries and documents as sets of contextualized 
token vectors instead of a single embedding. This representational richness comes at a 
query-time cost: the MaxSim scoring operator compares every query token with every 
document token, so exhaustive scoring scales with the product of query length, document 
length, corpus size, and embedding dimension. Existing systems work, most prominently 
ColBERTv2 \citep{santhanam2022colbertv2} and PLAID \citep{santhanam2022plaid}, address 
this cost with compression and approximate candidate pruning.

An orthogonal line of database research reduces the cost of high-dimensional similarity 
search by decomposing vectors vertically, that is, by dimension. BOND \citep{devries2002bond} 
scans dimensions incrementally and uses partial distances together with residual bounds 
to prune objects that cannot enter the top-k. PDX \citep{kuffo2025pdx} revisits this dimension-major 
layout for modern hardware and combines it with dimension-pruning search algorithms, including a 
PDX-BOND variant. These techniques were designed for single-vector nearest-neighbor search. Whether 
they transfer to the multi-vector MaxSim operator, where the scored object is a variable-length 
token matrix and the score is a sum of token-wise maxima, is an open engineering question that we 
set out to answer.

Our primary research question is whether BOND-style, exact-safe dimension pruning can accelerate 
exact ColBERT MaxSim search over a strong compiled exhaustive implementation (PRQ). To support our 
investigation, we also ask three secondary questions:
\begin{itemize}[noitemsep]
  \item \textbf{RQ1.} Why do the BOND residual bounds become useful, or remain loose, on normalized 
  ColBERT token embeddings?
  \item \textbf{RQ2.} Can a flat token-vector index generate document candidates with high exact-MaxSim recall?
  \item \textbf{RQ3.} At matched rerank budgets, how do PDX-IVF and FAISS-IVF trade exact-ranking 
  recovery for online latency?
  \item \textbf{RQ4.} Under the same machine, query split, and complete online
    timing boundary, how does a validation-selected CPU PLAID configuration
    compare with the exact and IVF pipelines?
\end{itemize}
We explicitly do not claim inverted-file (IVF) candidate generation as a novel contribution. IVF 
is standard approximate nearest-neighbor practice; in this project it serves as a practical comparison 
architecture and as a fallback that we adopted while testing the BOND hypothesis.

\subsection{Contributions}

The report makes the following evidence-based contributions:

\begin{enumerate}[noitemsep]
  \item A packed, flattened token-index representation of variable-length
    ColBERT embeddings, with a token-to-document mapping that supports both
    vertically decomposed (PDX) and conventional (FAISS) indexes over the same
    data (\cref{sec:architecture}).
  \item A controlled benchmark protocol with a single machine and process,
    matched work budgets, complete online timing boundaries, pinned threads,
    interleaved repetitions, and audited release artifacts
    (\cref{sec:methodology}). An internal measurement audit that led to the
    withdrawal of earlier, uncontrolled speedup claims is summarized as part
    of the methodology.
  \item A controlled evaluation of IVF candidate generation followed by exact
    MaxSim reranking on SciFact and NFCorpus~\citep{thakur2021beir}, showing
    identical FAISS-IVF and PDX-IVF recovery curves and no consistent
    engine-specific latency advantage (\cref{sec:results-ivf}).
  \item A controlled CPU PLAID comparison with predeclared validation-based
    parameter selection, held-out evaluation, and a separately labelled
    post-hoc full-score sensitivity check. The result documents a strong
    validation--test shift and prevents the earlier unmatched PLAID comparison
    from being reused as performance evidence (\cref{sec:results-plaid}).
  \item An independently implemented, provably exact-safe BOND-MaxSim kernel
    and a matched-precision comparison against a compiled exhaustive MaxSim
    baseline. The central result is negative: on raw ColBERT dimensions the
    kernel evaluates 94--95\% of the exhaustive component products and is
    roughly ten times slower than the exhaustive baseline
    (\cref{sec:results-bond}).
  \item A mechanism analysis using oracle pruning thresholds and an offline
    PCA rotation, showing that dimension ordering substantially improves
    pruning (down to 61.5\% of component products) yet still does not
    overcome the kernel's bound-maintenance overhead
    (\cref{sec:results-pca}).
\end{enumerate}

Throughout the report, approximate IVF results and exact-safe BOND results 
are kept strictly separate, and we do not claim any end-to-end retrieval-system speedup. 
All latency statements refer to a declared online query-processing boundary 
and to arms measured in the same interleaved run.

% \subsection{Report organization}

% \Cref{sec:background} reviews late interaction, MaxSim, PDX, BOND, IVF, and
% related adaptive-distance work. \Cref{sec:architecture} describes the data
% representation and the system components. \Cref{sec:methodology} defines the
% controlled benchmark protocol. \Cref{sec:results-ivf,sec:results-plaid,sec:results-bond,sec:results-pca}
% present the IVF, PLAID, BOND, and PCA results. \Cref{sec:discussion} interprets the
% evidence, \cref{sec:limitations} discusses threats to validity,
% \cref{sec:reproducibility} documents reproducibility measures, and
% \cref{sec:conclusion} concludes.
